% =====================================================================
%  PARTIE I : exemple d'exercice avec du code Python
% =====================================================================
\chapter{Prise en main de Python et structures de contrôle}

\section{Énoncé de l'exercice}

\begin{remarque}[Consigne]
Écrire un programme Python qui calcule la moyenne d'une liste de notes, puis
affiche la mention correspondante en utilisant la structure conditionnelle
\codeinline{if / elif / else}.
\end{remarque}

Le programme doit demander à l'utilisateur de saisir les notes, calculer la
moyenne et afficher un message indiquant si l'étudiant est admis.

\section{Démarche et implémentation}

Nous définissons d'abord une fonction \codeinline{moyenne} qui prend une liste
en paramètre et renvoie sa moyenne. La fonction \codeinline{mention} associe
ensuite une mention à la moyenne obtenue.

\begin{code}[language=Python]{notes.py}
# Calcul d'une moyenne et affichage de la mention associee

def moyenne(notes):
    """Retourne la moyenne d'une liste de notes."""
    if not notes:
        return 0.0
    return sum(notes) / len(notes)


def mention(moy):
    """Retourne la mention correspondant a une moyenne."""
    if moy >= 16:
        return "Tres bien"
    elif moy >= 14:
        return "Bien"
    elif moy >= 12:
        return "Assez bien"
    elif moy >= 10:
        return "Admis"
    else:
        return "Ajourne"


if __name__ == "__main__":
    notes = [15.5, 12.0, 17.0, 9.5]
    moy = moyenne(notes)
    print(f"Moyenne : {moy:.2f}")
    print(f"Mention : {mention(moy)}")
\end{code}

\section{Résultats et interprétation}

L'exécution du programme dans le terminal produit la sortie suivante.

\begin{terminal}
$ python notes.py
Moyenne : 13.50
Mention : Assez bien
\end{terminal}

La moyenne calculée vaut \textbf{13{,}50}, ce qui correspond à la mention
«~Assez bien~». Le test de la fonction \codeinline{mention} avec différentes
valeurs permet de vérifier que chaque branche conditionnelle renvoie la
mention attendue (voir Tableau~\ref{tab:mentions}).

\begin{table}[h]
  \centering
  \begin{tabularx}{0.75\linewidth}{@{}Xc@{}}
    \toprule
    \textbf{Intervalle de la moyenne} & \textbf{Mention} \\
    \midrule
    $[16~;~20]$   & Très bien  \\
    $[14~;~16[$   & Bien       \\
    $[12~;~14[$   & Assez bien \\
    $[10~;~12[$   & Admis      \\
    $[0~;~10[$    & Ajourné    \\
    \bottomrule
  \end{tabularx}
  \caption{Correspondance entre moyenne et mention.}
  \label{tab:mentions}
\end{table}

\begin{remarque}[Interprétation]
Le programme respecte la consigne. On remarque que l'absence de notes est
gérée par le cas de base \codeinline{not notes} afin d'éviter une division par
zéro. Toutes les valeurs limites ont été testées.
\end{remarque}
